problem:  
Let \( M^{n-1} \subset S^n \) be a closed, embedded minimal hypersurface of the unit sphere.  
Let \(|A|^2\) denote the squared norm of its second fundamental form and define  
\[
\operatorname{Var}_M(|A|^2) = \frac{1}{\operatorname{Vol}(M)} \int_M \big(|A|^2 - \overline{|A|^2}\big)^2\, dV,
\qquad 
\overline{|A|^2} = \frac{1}{\operatorname{Vol}(M)}\int_M |A|^2\, dV.
\]
Let \( \lambda_1, \lambda_2 \) be the first two nonzero eigenvalues of the Laplace–Beltrami operator on \(M\).  

**Open Problem (Quantitative Spectral Rigidity for Minimal Hypersurfaces in \(S^n\)):**  
Assuming Yau’s conjecture \( \lambda_1(M) = n - 1 \) holds for all minimal hypersurfaces \(M \subset S^n\), determine whether there exists a universal positive dimensional constant \( \kappa_n > 0 \) such that  
\[
\lambda_2(M) - \lambda_1(M) \ge \kappa_n\, \operatorname{Var}_M(|A|^2)^{1/2},
\]
and that equality (or asymptotic equality as \( \operatorname{Var}_M(|A|^2)\to 0\)) holds if and only if \( M \) is an isoparametric minimal hypersurface (e.g. equator, Clifford torus, or other homogeneous examples).  

Equivalently, can the *spectral gap* be quantitatively bounded below by the deviation of \(|A|^2\) from constancy, indicating that uniformity of extrinsic curvature forces spectral clustering, and its non-uniformity enforces a definite spectral separation?  

Why is it a "good" problem:  
- **Mathematically consistent and compatible with Yau’s conjecture:** The problem avoids contradictions with the strong belief that \( \lambda_1 = n-1 \) by focusing instead on higher spectral data and on quantitative deviation from isoparametricity.  
- **Precise and falsifiable:** The formulation is no longer vacuous—it specifies an explicit dependence (square-root growth) and a sharp condition for equality, making it a concrete, testable conjecture within analytic geometry.  
- **Bridges geometry and spectrum:** It asks whether *spectral gaps* reflect variation of *extrinsic curvature*, connecting the spectrum of the Laplacian (an intrinsic analytic invariant) to quantitative measures of extrinsic geometric irregularity.  
- **Heuristics from Simons identity:** The Laplacian of \(|A|^2\) on a minimal hypersurface satisfies a second-order differential relation (Simons identity) involving \(|\nabla A|^2\) and \(|A|^4\); integrating these identities naturally leads to expressions involving variance of \(|A|^2\). Thus, a quantitative dependence between curvature variance and spectrum is a plausible analytic direction.  
- **Depth and breadth:** A proof or even partial estimate would deepen the understanding of how geometric uniformity controls high-order spectral behavior and could open pathways to new rigidity or stability results for minimal submanifolds.  
- **Simplicity and universality:** The conjecture is compactly stated using familiar quantities (\(\lambda_i\), \(|A|^2\)), yet its resolution would require substantial advances combining spectral theory, curvature analysis, and geometric PDE.  

Hence, it is a *good* problem: simple to state, deeply rooted in core structures of minimal submanifold and spectral geometry, logically coherent, and societally stimulating for geometric analysts seeking quantitative links between intrinsic and extrinsic phenomena.